\documentclass[10pt,a4paper]{article}
\usepackage[margin=2cm]{geometry}
\usepackage[T1]{fontenc}
\usepackage{lmodern,booktabs,tabularx}
\usepackage[hidelinks]{hyperref}
\usepackage{microtype}
\setlength{\parindent}{0pt}
\setlength{\parskip}{5pt}
\title{MLIP training environments: GPU checks and foundation-model readiness}
\author{MLIP-nucleation project}
\date{17 October 2026}
\begin{document}
\maketitle

This report preserves the initial DeepMD environment observations. The subsequent
separate GRACE installation and successful A100 training check are recorded in
\path{reports/grace-installation-2026-10-17.tex}.

\section{Operational conclusion}
On the tested A100, the existing \texttt{lammps\_deepmd\_dpa4} conda environment
successfully trained small DPA4 and DPA4c models. Its name does not limit it to
LAMMPS: its DeePMD training executable was tested directly. The initial
\texttt{deepmd\_320\_a100} environment trained DPA4c but failed on DPA4 because
the Torch neighbor-list dependency was missing. No installed environment was modified.

The existing \texttt{deepmd\_2210\_k40} environment trained a conventional
\texttt{se\_e2\_a} model on K40m. This is a usable legacy software path;
it does not establish modern foundation-model support on that GPU.
GRACE/TensorPotential was not found in the inspected conda environments.

\section{Slurm runs and acceptance criteria}
All runs used account \texttt{imet-2026-2}, QoS \texttt{normal}, one GPU,
four CPUs and 16 GiB RAM. The initial two jobs had ten-minute limits.
After they finished in a combined 172 GPU-seconds, a five-minute follow-up
tested a different existing environment with the missing dependency.
The three jobs consumed 241 GPU-seconds in total (scheduler elapsed times,
including startup), below the initial twenty GPU-minute envelope.

\begin{center}
\begin{tabularx}{\linewidth}{@{}llllX@{}}
\toprule
Job & Partition/node & Elapsed & Outcome & Scope \\
\midrule
20950192 & v100/tesla-a101 & 1:27 & Partial & se\_e2\_a and DPA4c passed; DPA4 dependency failure \\
20950193 & debug/tesla47 & 1:25 & Passed & legacy se\_e2\_a training \\
20950195 & v100/tesla-a101 & 1:09 & Passed & se\_e2\_a, DPA4 and DPA4c training \\
\bottomrule
\end{tabularx}
\end{center}

Each successful case completed 20 optimizer steps, produced a checkpoint,
and had a finite training/validation learning curve. A separate GPU probe
performed forward and gradient operations on an explicitly selected CUDA/GPU
device; DeepMD logs independently identified the training device as CUDA:0 or GPU:0.
Scheduler completion alone was not used as the acceptance criterion.

The inputs are deliberately small synthetic harmonic-pair configurations:
four Al-labelled atoms, eight training and four independently generated
validation configurations, seed 20261017. Labels are analytic, not DFT and
not foundation-model predictions. Units are eV and eV/angstrom; the saved
virial convention is $\sum_i\mathbf r_i\otimes\mathbf F_i$, in eV and
row-major tensor order. Virial loss was disabled in these training probes.
This test neither selects Al as the research system nor measures scientific
accuracy, production throughput, memory scaling or model suitability.
Small-run timings across the two software stacks are not a GPU performance comparison.

\section{Observed software and failure diagnosis}
The A100 was an NVIDIA A100-PCIE-40GB with driver 580.126.20 and compute
capability 8.0. Both tested A100 environments contain DeePMD 3.2.0.
The first used PyTorch 2.12.0/CUDA 13.0; the successful follow-up used
PyTorch 2.11.0+cu130. Their installed Torch architecture lists include
\texttt{sm\_80}, but not K40m's \texttt{sm\_35}.

In \texttt{deepmd\_320\_a100}, ordinary \texttt{vesin} 0.6.1 is present,
but \texttt{vesin-torch} is absent. The DPA4 traceback reports that neither
neighbor-list implementation is importable. Inspection of the installed
selection function identifies \texttt{vesin-torch} or
\texttt{nvalchemi-toolkit-ops} as the relevant alternatives.
The follow-up environment already contains \texttt{vesin-torch} 0.6.1 and
passes the identical DPA4 input. This supports using that existing environment
for initial DPA4/DPA4c work. Adding \texttt{vesin-torch} to a separately
versioned clone of the first environment is an alternative; an in-place
upgrade was not performed. DPA4c used the available dense neighbor-graph
fallback, so the optional NV acceleration path remains untested.

The K40m exposed 11441 MiB of GPU memory and driver 440.33.01.
Its environment contains DeePMD 2.2.10 and TensorFlow 2.9.0 built for
CUDA 10.2/cuDNN 7, explicitly including \texttt{sm\_35}. Both its GPU
gradient probe and DeepMD training passed. Preserve this environment for
compatible legacy students; the modern A100 Torch binaries cannot be used
on K40m. This is a software compatibility observation, not a claim that all
possible older or custom builds have been evaluated.

\section{What remains for foundation-model fine-tuning}
\begin{enumerate}
\item \textbf{GRACE installation.} Create a separate environment, initially
Python 3.11, TensorPotential 0.6.1, TensorFlow 2.20 and matching legacy Keras.
Published package metadata requires TensorFlow $\leq2.20$ on Linux,
including CUDA dependencies. Existing DeepMD A100 TensorFlow is 2.21, so
reusing it would introduce a dependency conflict.
\path{environments/grace-a100.proposed.yml} is an installation proposal,
not an installed or resolved lockfile. Resolve its dependencies, capture
an explicit lock and verify actual GRACE training on A100 before use.

\item \textbf{Trainable checkpoint selection.} No checkpoint was selected
or downloaded by these tests, and standard DeepMD/GRACE cache directories
were absent. Choose a checkpoint separately from the architecture, recording
source revision, file hash, license, chemical coverage and pretraining label
settings. GRACE fine-tuning needs checkpoint files and model definition,
not only an inference SavedModel. DeepMD exposes a pretrained-model
download command; inspect the installed version's model list rather than
assuming every architecture preset is pretrained weights.

\item \textbf{Real fine-tuning and deployment checks.} Test loading the
chosen weights, energy-reference alignment, parameter updates, save/restart,
and export/inference consistency. Those checks, LoRA, multi-GPU training,
DPA4c compression and LAMMPS execution were not exercised here. Optional
\texttt{nvalchemi-toolkit-ops} acceleration should be evaluated separately.
Use the independently validated DFT protocol and grouped data splits for
scientific work, preserving DFT and teacher labels separately.
\end{enumerate}

\section{Persistence and extension}
The candidate architecture registry is \path{configs/model_backends.json}.
It separates architecture, execution backend, candidate role and a future
checkpoint record, and lists the operations required of a future shared
adapter. Such an adapter has not yet been implemented. Additional architectures
can be added without changing the research workflow or silently replacing
the selected checkpoint.

Run records are under \path{runs/mlip-smoke-20261017-a100-01/},
\path{runs/mlip-smoke-20261017-k40m-01/} and
\path{runs/mlip-smoke-20261017-a100-vesin-01/}.
They retain submission intents, job IDs, input script hashes, generated-data
hashes, exact training inputs, package snapshots, scheduler state, learning
curves and checkpoint hashes. Collected records are checked against remote
SHA-256 values. Test checkpoints remain in their remote run directories under
\path{/home6/u3129/mlip-research/runs/}; they are not accepted scientific models
or durable model releases. The runtime source is preserved by Git commit
\texttt{6bfe80f}; individual manifest hashes identify the exact submitted bytes.

\section{Primary sources consulted}
Accessed 17 October 2026. Cluster observations above come from the saved run records.
\begin{itemize}
\item DPA4 backend and fine-tuning guidance:
\url{https://docs.deepmodeling.com/projects/deepmd/en/stable/model/dpa4.html}.
\item DPA4c backend and compression requirements:
\url{https://docs.deepmodeling.com/projects/deepmd/en/stable/model/dpa4c.html}.
\item GRACE installation and legacy Keras:
\url{https://gracemaker.readthedocs.io/en/latest/gracemaker/install/}.
\item Published TensorPotential 0.6.1 metadata:
\url{https://pypi.org/pypi/tensorpotential/0.6.1/json}.
\item GRACE checkpoint requirements:
\url{https://gracemaker.readthedocs.io/en/latest/gracemaker/foundation/}.
\item DeepMD pretrained-model selection:
\url{https://docs.deepmodeling.com/projects/deepmd/en/stable/model/pretrained.html}.
\end{itemize}
\end{document}
